\documentclass[10pt,a4paper]{amsart}
\usepackage[margin=19mm]{geometry}
\usepackage{amsmath,amssymb,mathtools}
\usepackage[scr=rsfs,cal=euler]{mathalfa}
\usepackage{xfrac}
\usepackage[unicode,hidelinks,bookmarks=false]{hyperref}
\newcommand{\ie}{i.e.\ }
\newcommand{\cf}{cf.\ }
\newcommand{\resp}{respectively\ }
\newcommand{\Subsection}{\subsection}
\providecommand{\nobreakdash}{\nobreak\hbox{-}}
\setlength{\emergencystretch}{2em}
\multlinegap0pt
\usepackage{bm}
\usepackage{bbm}
\usepackage{dsfont}
\usepackage{xspace}
\usepackage{cancel}
\usepackage{mathtools}
\usepackage[shortlabels]{enumitem}
\usepackage{amsthm}
\makeatletter
\@ifundefined{theorem}{\newtheorem{theorem}{Theorem}}{}
\makeatother
\newcommand{\R}{\mathbb{R}}
\newcommand{\bmat}[1]{\begin{bmatrix} #1 \end{bmatrix}}
\newcommand{\ip}[2]{\left\langle #1,#2\right\rangle}
\newcommand{\mbf}[1]{\mathbf{#1}}
\newcommand{\mcl}[1]{\mathcal{#1}}
\newcommand{\norm}[1]{\lVert#1\rVert}
\title[Impulse-to-Peak-Output Norm Optimal State-Feedback Control o]{Impulse-to-Peak-Output Norm Optimal State-Feedback Control of Linear PDEs}
\author{Tristan Thomas, Sachin Shivakumar, Javad Mohammadpour Velni}
\date{Source: arXiv:2604.03399v2 (CC BY 4.0)}
\begin{document}
\maketitle

\noindent\emph{Source and licence.} Impulse-to-Peak-Output Norm Optimal State-Feedback Control of Linear PDEs. Version arXiv:2604.03399v2 (CC BY 4.0), distributed under the Creative Commons Attribution 4.0 International licence, as recorded in the arXiv licence metadata for this exact version and shown on the abstract page https://arxiv.org/abs/2604.03399v2. Full rights notice: licenses/arxiv-2604.03399-CC-BY-4.0.txt.\par\smallskip

\noindent\textit{Editorial scope.} Counted unit: the Theorem whose source label is \texttt{thm:Normnoncoercive} in the article (source0.tex lines 468--484, source label \texttt{thm:Normnoncoercive}), delivered together with the article's own complete original proof of it (source0.tex lines 486--522, one \texttt{proof} environment of 37 lines). No mathematics is written, repaired or re-derived: statement and proof are the article's own text line for line. Required context retained uncounted: eqn:PIEaux (lines 313--315); thm:Normcoercive (lines 437--445). Every definition, display and label of those lines is retained unchanged. Omitted from this package and not redistributed: 1--312, 316--436, 446--467, 485--485, 523--824. Nothing inside a retained excerpt is omitted or altered: every retained body is delivered exactly as the article's lines are written, so no omission receipt of any kind accompanies this package. Citations: the retained excerpts carry no citation command at all, so no bibliography entry of the article is transcribed and no reference is redistributed. Reference adaptation: none was needed and none was made; every reference inside the delivered unit resolves inside it, so no unresolved reference or citation is produced. Printed slips kept literal: no printed word, symbol, hypothesis or number of the counted statement or of its proof was altered. Layout and class: the article's own class is \texttt{ieeeconf} with packages including comment, booktabs, cite, mathtools, amsmath, amssymb, amsfonts, algorithmic, graphicx, textcomp, makecell, xcolor, amsthm, microtype; the wrapper is the lane's pinned \texttt{amsart} editorial wrapper at 10pt on A4, which restates the theorem-like environments the retained text needs and adds the article's own macro definitions that the retained text needs (\texttt{\textbackslash R}, \texttt{\textbackslash bmat}, \texttt{\textbackslash ip}, \texttt{\textbackslash mbf}, \texttt{\textbackslash mcl}, \texttt{\textbackslash norm}), with extra wrapper packages (bm, bbm, dsfont, xspace, cancel, mathtools, enumitem). The delivered numbering is the unit's own. Assets: no retained span contains a figure environment, a graphics-inclusion command, a PDF-page inclusion, a table or a TikZ picture; no image file is copied into this package, the package declares no asset dependency and no image file is redistributed with it. Licence: version arXiv:2604.03399v2 of this article is distributed under the Creative Commons Attribution 4.0 International licence, as recorded in the arXiv licence metadata for this exact version and shown on the abstract page https://arxiv.org/abs/2604.03399v2. A full rights notice is delivered at licenses/arxiv-2604.03399-CC-BY-4.0.txt. Characters: every retained excerpt is pure ASCII, so the delivered engine, XeLaTeX with its default Latin Modern text font, reports no missing character.
\begin{align}\label{eqn:PIEaux}
\partial_t (\mcl T \mbf x(t))&= \mcl A \mbf x(t),\quad \mcl T \mbf{x}(0) =\mcl{B} v,\quad z(t)=\mcl C \mbf x(t),
\end{align}

\begin{theorem}\label{thm:Normcoercive}
    Given 4-PI operators ${\mathcal{T},\mathcal{A},\mathcal{B},\mathcal{C}}$, let $\{\mbf x, z\}$ satisfy a PIE of the form
    Eq. \eqref{eqn:PIEaux}. Suppose there exists a $\gamma>0$ and a 4-PI operator $\mcl Q \succeq 0$ such that one of the following two conditions hold:
    \begin{enumerate}[leftmargin=*]
        \item[(i)] $\mathcal{B}^*\mcl Q\mathcal{B}\preceq I$, $\mathcal{A}^*\mcl Q\mathcal{T}+\mathcal{T}^*\mcl Q\mathcal{A} \preceq 0$, $\frac{1}{\gamma^2}\mathcal{C}^*\mathcal{C} \preceq \mathcal{T}^*\mcl Q\mathcal{T}$.
        \item[(ii)] $\mathcal{C}\mcl Q\mathcal{C}^*\preceq I$, $\mathcal{A}\mcl Q\mathcal{T}^*+\mathcal{T}\mcl Q\mathcal{A}^* \preceq 0$, $\frac{1}{\gamma^2}\mathcal{B}\mathcal{B}^* \preceq \mathcal{T}\mcl Q\mathcal{T}^*$.
    \end{enumerate} 
    Then, $\norm{G(\mcl T, \mcl A,\mcl B,\mcl C)}_{ip} \le \gamma$.
\end{theorem}

\begin{theorem}\label{thm:Normnoncoercive}
    Given 4-PI operators ${\mathcal{T}, \mathcal{A},\mathcal{B},\mathcal{C}}$, let $\{\mbf x, z\}$ satisfy a PIE of the form Eq. (\ref{eqn:PIEaux}).
    Suppose there exists a $\gamma > 0$ and a 4-PI operator $\mcl Q$ such that $\mcl T^*\mcl Q =\mcl Q^* \mcl T \succeq 0$ and one of the following conditions hold:
    \begin{enumerate}
        \item[(i)] $\mcl A^*\mcl Q+\mcl Q^*\mcl A \!\preceq\! 0$, $ \begin{bmatrix}
            \gamma^2I & \mcl C\\ \mcl C^* & \mcl Q^* \mcl  T 
        \end{bmatrix}\!\succeq\! 0$, $ \begin{bmatrix}
            \mcl T^*\mcl Q & \mcl Q^* \mcl B\\ \mcl B^* \mcl Q & I
        \end{bmatrix}\!\succeq\! 0$ 
        \item[(ii)] $\mcl A\mcl Q+\mcl Q^*\mcl A^* \!\preceq\! 0$, $ \begin{bmatrix}
            \gamma^2I \!\!& \mcl B^*\\ \mcl B \!\!& \mcl Q^* \mcl  T^*
        \end{bmatrix}\!\succeq\! 0$, $ \begin{bmatrix}
            \mcl T\mcl Q & \mcl Q^* \mcl C^*\\ \mcl C \mcl Q & I
        \end{bmatrix}\!\succeq\! 0$ 
    \end{enumerate} 
    Then, $\norm{G(\mcl T, \mcl A,\mcl B,\mcl C)}_{ip} \le \gamma$.
\end{theorem}

\begin{proof}
We give a proof outline, since it is similar to the proof of Thm. \ref{thm:Normcoercive}.
Let $\{\mbf x, z\}$ satisfy the PIE Eq. \eqref{eqn:PIEaux} with initial conditions $\mcl T\mbf x(0)=\mcl Bv$. Next, define a storage function as $V(\mbf x) = \ip{\mcl T\mbf x}{\mcl Q\mbf x}$. Then, if $\mcl Q$ satisfies the constraints in \textit{(i)}, we have
$\dot{V}(\mbf x(t))\le 0$ and $\norm{z(t)}^2\le \gamma^2V(t)$ for all $t>0$.
% \begin{align*}
%     V(\varphi) = \ip{\varphi}{\mcl T^* \mcl Q \varphi} = \ip{\varphi}{\mcl Q^* \mcl T \varphi}\\
%     \dot{V}(\varphi) = \ip{\varphi}{\mcl Q^* \partial_t \mcl T \varphi} + \ip{\partial_t \varphi}{\mcl Q^* \mcl T \varphi}\\
%     = \ip{\varphi}{\mcl Q^* \partial_t \mcl T \varphi} + \ip{\mcl Q^*\mcl T\partial_t \varphi}{\varphi}\\
%     =\ip{\varphi}{\mcl Q^* \mcl A \varphi} + \ip{\mcl Q^*\mcl A \varphi}{\varphi}\\ 
%     = \ip{\varphi}{\mcl Q^* \mcl A \varphi}+\ip{\varphi}{\mcl A^* \mcl Q \varphi}\\
%     = \ip{\varphi}{(\mcl Q^* \mcl A+\mcl A^*\mcl Q)\varphi}
% \end{align*}
% Using assumption 1, the above storage function, with negative semi-definite time derivative, is non-increasing in time. 
% \begin{comment}
       
    
%         \item Take the Lyapunov function $V(\varphi) = \ip{\varphi}{\mcl R \varphi} = \ip{\varphi}{\mcl Q^* \mcl T \varphi} = \ip{\mcl Q \varphi}{\mcl T \varphi}$. We can rewrite the first part of the inner product such that the total expression gives $\ip {\mcl Q \mcl T^{-1} \mcl T \varphi}{\mcl T \varphi}= \ip{\mcl Q \mcl T^{-1} x}{x}$, where in this proof $x$ represents the PDE state. Taking the time derivative of the storage function, $D_tV = \ip{\mcl Q \mcl T^{-1} \partial_t x}{x}+ \ip{\mcl Q \mcl T^{-1} x}{\partial_tx} = \ip{\mcl Q \mcl T^{-1} \mcl A \varphi}{x}+\ip{\mcl Q \mcl T ^{-1}x}{\mcl A \varphi}$. The second inner product term can utilize $x = \mcl T \varphi$ to map $\mcl T ^{-1} x$ back to $\varphi$. Thus the expression becomes $D_tV = \ip{\varphi}{\mcl A \mcl T^{-*}\mcl Q ^*x}+ \ip{\varphi}{\mcl Q ^* \mcl A \varphi}$. Because $\mcl Q^*\mcl T = \mcl T^* \mcl Q \implies \mcl T^{-*}\mcl Q^*\mcl T = \mcl T^{-*}\mcl T^*Q = Q.$ Thus, $D_tV = \ip{\varphi}{\mcl A^* \mcl T^{-*}\mcl Q^* \mcl T \varphi}+\ip{\varphi}{\mcl Q^* \mcl A \mcl T} = \ip{\varphi}{\mcl A^* \mcl Q  \varphi}+\ip{\varphi}{\mcl Q^* \mcl A \mcl T} = \ip{\varphi}{(\mcl A^* \mcl Q + \mcl Q^*\mcl A)\varphi}$. Thus, if $\mcl A^* \mcl Q + \mcl Q^*\mcl A \preceq 0$, the overall storage function time derivative will be nonpositive.
% \end{comment}

        % Next, $\norm{z(t)}^2 = \ip{z(t)}{z(t)} = \ip{\mcl C \varphi}{\mcl C \varphi} = \ip{\varphi}{\mcl C^*\mcl C\varphi}.$ If the second condition is met, by Schur Complement, $\mcl C^* \mcl C \preceq \mcl Q^* \mcl T$, then $\ip{\varphi}{\mcl C^* \mcl C\varphi}\preceq \gamma^2 \ip{\varphi}{\mcl Q^* \mcl T\varphi}= \gamma^2 \ip{\mcl Q \varphi }{\mcl T \varphi} = \gamma^2 V(t)$
        % Finally, we show that the storage function at time 0 is bounded by the initial storage function. $V = \ip{\mcl Q\varphi}{\mcl T\varphi}$. Then using $\mcl T \varphi = \mcl B v$, 
        Lastly, for any $v\in \R^{n_w}$ and $\mcl T\mbf x(0)=\mcl Bv$, we have
        \begin{align*}
            0\le &\ip{\bmat{\mbf x(0)\\-v}}{ \bmat{\mcl T^*\mcl Q&\mcl Q^*\mcl B\\\mcl B^*\mcl Q&I} \bmat{\mbf x(0)\\-v}} \\
            &= \ip{\mbf x(0)}{\mcl T^*\mcl Q\mbf x(0)} - 2\ip{\mbf x(0)}{\mcl Q^*\mcl Bv} + \norm{v}^2 \\
            &=\ip{\mbf x(0)}{\mcl Q^*\mcl T\mbf x(0)} - 2\ip{\mbf x(0)}{\mcl Q^*\mcl T\mbf x(0)} + \norm{v}^2 \\
            &= \norm{v}^2-V(\mbf x(0)).
        \end{align*}
        Since $\norm{z(t)}^2\le \gamma^2V(\mbf x(t))\le \gamma^2V(\mbf x(0))\le \gamma^2\norm{v}^2$, we have $\norm{G(\mcl T, \mcl A,\mcl B,\mcl C)}_{ip} \le \gamma$.
        Likewise, satisfaction of \textit{(ii)} implies the norm bound on the dual and the primal PIE response.
       % \begin{align*}
        % &V(\mbf x(0)) \!= \!\ip{\mcl Q \mbf x(0)}{\mcl B v} \!= \!\ip{\mbf x(0)}{\mcl Q^* \mcl B v} \!= \!\ip{\mcl T^{-1}\mcl T\mbf x(0)}{\mcl Q^* \mcl B v} \\
        % &= \ip{\mcl T ^{-1}\mcl B v}{\mcl Q ^ * \mcl Bv}= \ip{v}{\mcl B ^* \mcl T ^{-*} \mcl Q^* \mcl B v}
        % \end{align*}
        % Using assumption from the theorem,$ \ip{v}{\mcl B ^* \mcl T ^{-*} \mcl Q^* \mcl B v} \preceq \ip{v}{Iv} = \ip{v}{v} = \norm{v(t)}^2$. Thus, $V(0) \leq \norm{v^2}$. 
        % The final argument chains together similarly to the previous theorem. The last condition tells you that the initial storage function is bounded, the first condition tells you that the storage function decreases, and the second condition tells you the output is also bounded by the storage function scaled. We can chain everything together as $\norm{z(t)}^2\leq \gamma^2V(t) \leq \gamma^2 V(0)\leq \gamma^2 \norm{v}^2 $. Taking initial condition with norm 1, then $\norm{z(t)}^2 \leq \gamma^2.$ Finally $\sup_{\norm{v}= 1}\sup_{t\geq 0}\norm{z(t)}^2\leq \gamma^2$. 
\end{proof}

\end{document}
